% GENERATED FILE. Edit canonical lessons, then run npm run generate:books.
% canonical-source-hash: fnv1a64:2367c375e46d9b5a
% canonical-lessons: FR-C351-reagir, FR-C351-reaction, FR-C351-pretexte, FR-C351-justifier, FR-C351-justification

\chapter{Comment réagir}
\label{ch:fr-comment-reagir}
\hlchaptermodality{351}

\begin{chapteropening}
\textbf{By the end of this chapter:} \emph{I can describe how someone reacted, tell a real reason from a pretext, and justify a decision.}

Complete the last lesson of chapter 351: I can describe how someone reacted, tell a real reason from a pretext, and justify a decision.
\end{chapteropening}

\section[réagir]{réagir — to react}

\label{lesson:FR-C351-reagir}



\begin{quote}
Before the new one: say the French for to arouse, then the French for the motive.
\end{quote}

\subsection*{You'll want to know: réagir}
\textbf{réagir} — \textquotedblleft{}to react\textquotedblright{}.

\begin{quote}
Comment as-tu réagi quand tu as appris la nouvelle ?
\end{quote}

\begin{quote}
\emph{How did you react when you heard the news?}
\end{quote}

\emph{Réagir à quelque chose}. A \emph{-ir} verb like \emph{finir}.

\begin{cousinweb}
\emph{réagir} is built from \emph{ré-} \textquotedblleft{}back\textquotedblright{} and \emph{agir} \textquotedblleft{}to act\textquotedblright{}, which comes from Latin \emph{agere} \textquotedblleft{}to do, to drive\textquotedblright{} (Lewis \& Short).
\end{cousinweb}

\begin{grammarlens}[title={What you've built}]
One more word for telling, explaining and describing at B1.
\end{grammarlens}

\subsection*{Your turn}
\begin{itemize}
\raggedright
  \item \emph{Say it:} \emph{réagir}
  \item \emph{Say it:} \emph{réagir}, once more
  \item \emph{Say it:} a sentence of your own with \emph{réagir}
  \item \emph{Recall:} say the French for to arouse, then the French for the motive, then say \emph{réagir} again
\end{itemize}

\begin{tcolorbox}[breakable,colback=teal!4,colframe=teal!35!black,title={Before you move on}]
What does réagir mean? (\textquotedblleft{}To react\textquotedblright{}.) Say it once more.
\end{tcolorbox}
\section[la réaction]{la réaction — the reaction}

\label{lesson:FR-C351-reaction}



\begin{quote}
Before the new one: say the French for the motive, then the French for to react.
\end{quote}

\subsection*{You'll want to know: la réaction}
\textbf{la réaction} — \textquotedblleft{}the reaction\textquotedblright{}.

\begin{quote}
Sa réaction m'a surpris.
\end{quote}

\begin{quote}
\emph{Her reaction surprised me.}
\end{quote}

Feminine: \emph{une réaction}.

\begin{cousinweb}
\emph{réaction} comes from Medieval Latin \emph{reactiō} (EtymDB), built from \emph{re-} \textquotedblleft{}back\textquotedblright{} and \emph{āctiō} \textquotedblleft{}a doing\textquotedblright{}, from \emph{agere} \textquotedblleft{}to do\textquotedblright{} (Lewis \& Short).
\end{cousinweb}

\begin{grammarlens}[title={What you've built}]
One more word for telling, explaining and describing at B1.
\end{grammarlens}

\subsection*{Your turn}
\begin{itemize}
\raggedright
  \item \emph{Say it:} \emph{la réaction}
  \item \emph{Say it:} \emph{la réaction}, once more
  \item \emph{Say it:} a sentence of your own with \emph{la réaction}
  \item \emph{Recall:} say the French for the motive, then the French for to react, then say \emph{la réaction} again
\end{itemize}

\begin{tcolorbox}[breakable,colback=teal!4,colframe=teal!35!black,title={Before you move on}]
What does la réaction mean? (\textquotedblleft{}The reaction\textquotedblright{}.) Say it once more.
\end{tcolorbox}
\section[le prétexte]{le prétexte — the pretext}

\label{lesson:FR-C351-pretexte}



\begin{quote}
Before the new one: say the French for to react, then the French for the reaction.
\end{quote}

\subsection*{You'll want to know: le prétexte}
\textbf{le prétexte} — \textquotedblleft{}the pretext\textquotedblright{}.

\begin{quote}
Il a trouvé un prétexte pour partir plus tôt.
\end{quote}

\begin{quote}
\emph{He found a pretext to leave earlier.}
\end{quote}

\emph{Sous prétexte que} — on the pretext that.

\begin{cousinweb}
\emph{prétexte} is a learned borrowing from Latin \emph{praetextus} \textquotedblleft{}an outward show, a pretext\textquotedblright{} (EtymDB; Lewis \& Short), from \emph{praetexere} \textquotedblleft{}to weave in front, to cover\textquotedblright{}, built on \emph{texere} \textquotedblleft{}to weave\textquotedblright{}.
\end{cousinweb}

\begin{grammarlens}[title={What you've built}]
One more word for telling, explaining and describing at B1.
\end{grammarlens}

\subsection*{Your turn}
\begin{itemize}
\raggedright
  \item \emph{Say it:} \emph{le prétexte}
  \item \emph{Say it:} \emph{le prétexte}, once more
  \item \emph{Say it:} a sentence of your own with \emph{le prétexte}
  \item \emph{Recall:} say the French for to react, then the French for the reaction, then say \emph{le prétexte} again
\end{itemize}

\begin{tcolorbox}[breakable,colback=teal!4,colframe=teal!35!black,title={Before you move on}]
What does le prétexte mean? (\textquotedblleft{}The pretext\textquotedblright{}.) Say it once more.
\end{tcolorbox}
\section[justifier]{justifier — to justify}

\label{lesson:FR-C351-justifier}



\begin{quote}
Before the new one: say the French for the reaction, then the French for the pretext.
\end{quote}

\subsection*{You'll want to know: justifier}
\textbf{justifier} — \textquotedblleft{}to justify\textquotedblright{}.

\begin{quote}
Elle a justifié sa décision devant l'équipe.
\end{quote}

\begin{quote}
\emph{She justified her decision in front of the team.}
\end{quote}

\emph{Se justifier} — to explain yourself.

\begin{cousinweb}
\emph{justifier} comes from Late Latin \emph{iūstificāre} \textquotedblleft{}to justify\textquotedblright{} (EtymDB), built on \emph{iūstus} \textquotedblleft{}right, fair\textquotedblright{} (Lewis \& Short) and \emph{-ficāre} \textquotedblleft{}to make\textquotedblright{}.
\end{cousinweb}

\begin{grammarlens}[title={What you've built}]
One more word for telling, explaining and describing at B1.
\end{grammarlens}

\subsection*{Your turn}
\begin{itemize}
\raggedright
  \item \emph{Say it:} \emph{justifier}
  \item \emph{Say it:} \emph{justifier}, once more
  \item \emph{Say it:} a sentence of your own with \emph{justifier}
  \item \emph{Recall:} say the French for the reaction, then the French for the pretext, then say \emph{justifier} again
\end{itemize}

\begin{tcolorbox}[breakable,colback=teal!4,colframe=teal!35!black,title={Before you move on}]
What does justifier mean? (\textquotedblleft{}To justify\textquotedblright{}.) Say it once more.
\end{tcolorbox}
\section[la justification]{la justification — the justification}

\label{lesson:FR-C351-justification}



\begin{quote}
Before the new one: say the French for the pretext, then the French for to justify.
\end{quote}

\subsection*{You'll want to know: la justification}
\textbf{la justification} — \textquotedblleft{}the justification\textquotedblright{}.

\begin{quote}
Il n'y a aucune justification pour ce prix.
\end{quote}

\begin{quote}
\emph{There is no justification for this price.}
\end{quote}

Feminine: \emph{une justification}.

\begin{cousinweb}
\emph{justification} is a learned borrowing from Late Latin \emph{iūstificātiō} (EtymDB), from \emph{iūstificāre} \textquotedblleft{}to justify\textquotedblright{}, built on \emph{iūstus} \textquotedblleft{}right, fair\textquotedblright{} (Lewis \& Short).
\end{cousinweb}

\begin{grammarlens}[title={What you've built}]
One more word for telling, explaining and describing at B1.
\end{grammarlens}

\subsection*{Your turn}
\begin{itemize}
\raggedright
  \item \emph{Say it:} \emph{la justification}
  \item \emph{Say it:} \emph{la justification}, once more
  \item \emph{Say it:} a sentence of your own with \emph{la justification}
  \item \emph{Recall:} say the French for the pretext, then the French for to justify, then say \emph{la justification} again
\end{itemize}

\begin{tcolorbox}[breakable,colback=teal!4,colframe=teal!35!black,title={Before you move on}]
What does la justification mean? (\textquotedblleft{}The justification\textquotedblright{}.) Say it once more.
\end{tcolorbox}
